\begin{lemma}[The untruncated kernel integral of a smooth input off its support]
Fix $n \geq 1$ and $j \in \{1,\ldots,n\}$, let $\lambda$ be Lebesgue measure on
$\mathbb{R}^n$, let $\gamma_n$ be the Gaussian measure of \noderef{GaussMeasure}, and let
$K_{\gamma,j}$ be the normalized Gaussian Riesz kernel of \noderef{GaussRieszKernel}.  Let
$\varphi : \mathbb{R}^n \to \mathbb{R}$ be smooth with compact support, let
$z \in \mathbb{R}^n$, let $\sigma > 0$, and suppose that
\[
  \varphi(y) = 0 \qquad\text{whenever } |z-y| \leq \sigma .
\]
Write, for $s > 0$,
\[
  (A_sw)(z) = \int_{\mathbb{R}^n} w\bigl(e^{-s}z + \sqrt{1-e^{-2s}}\,\zeta\bigr)\,d\gamma_n(\zeta)
\]
for the Mehler average, and
\[
  \Phi_j^\varphi(z) = \frac{1}{\sqrt\pi}\int_0^\infty
      \frac{1}{\sqrt s}\,e^{-s}\,\frac{1}{\sqrt2}
      \bigl(A_s(\partial_j\varphi)\bigr)(z)\,ds ,
\]
which is the quantity appearing as the limit in clause (2) of \noderef{PVCoreSmooth} at the
point $z$.  Then the function $y \mapsto K_{\gamma,j}(z,y)\varphi(y)$ is Lebesgue integrable on
$\mathbb{R}^n$, and
\[
  \Phi_j^\varphi(z) = \int_{\mathbb{R}^n}K_{\gamma,j}(z,y)\,\varphi(y)\,dy .
\]

\emph{This is an off-diagonal identification and not a principal value.}  No limit in a
truncation radius is taken in the conclusion: the point of the hypothesis on $\varphi$ is that
on the range $0 < \varepsilon \leq \sigma$ the truncated integrals do not depend on
$\varepsilon$ at all, so the limit that \noderef{PVCoreSmooth} exhibits is attained, and is
attained by the integral over all of $\mathbb{R}^n$.  The reference paper licenses exactly this
much when it says that for compactly supported smooth inputs and a point off their support the
integrals may be interchanged directly, and immediately adds that this is an off-diagonal kernel
identification rather than a principal-value identification; the present statement stays on the
first side of that line.

\emph{The integrability assertion is content.}  The right-hand side is a Bochner integral, equal
to $0$ by convention on a non-integrable integrand; without the integrability assertion the
displayed identity could hold with the right-hand side conventional.  It is asserted for the
given $z$ with no exceptional set, because $\varphi$ is an honest function and the hypothesis is
pointwise.
\end{lemma}
\begin{proof}
Write $F(y) = K_{\gamma,j}(z,y)\varphi(y)$ and, for $\varepsilon > 0$,
$\Omega_\varepsilon = \{y : |z-y| > \varepsilon\}$, a measurable set (the preimage of an open
half-line under the continuous map $y \mapsto |z-y|$).

\emph{Step 1: for every $\varepsilon \in (0,\sigma]$, $F$ is integrable on $\mathbb{R}^n$ and
its integral is the truncated one.}
Fix such an $\varepsilon$.  Clause (1) of \noderef{PVCoreSmooth}, applied to $\varphi$ (smooth
with compact support), to $j$, to the point $z$ and to this $\varepsilon > 0$, states that $F$
is Lebesgue integrable on $\Omega_\varepsilon$.

Next, $F$ vanishes off $\Omega_\varepsilon$: if $y \notin \Omega_\varepsilon$ then
$|z-y| \leq \varepsilon \leq \sigma$, so $\varphi(y) = 0$ by hypothesis and hence $F(y) = 0$.
This holds at every such $y$, not merely almost every one.

Let $\mathbf 1_{\Omega_\varepsilon}F$ be the function equal to $F$ on $\Omega_\varepsilon$ and
to $0$ elsewhere.  Since $\Omega_\varepsilon$ is measurable,
$\mathbf 1_{\Omega_\varepsilon}F$ is Lebesgue integrable on $\mathbb{R}^n$ if and only if $F$ is
Lebesgue integrable on $\Omega_\varepsilon$ (shared context: the indicator criterion for
integrability on a measurable set), so it is.  And
$F = \mathbf 1_{\Omega_\varepsilon}F$ everywhere: the two agree on $\Omega_\varepsilon$ by
definition and both vanish off it by the previous paragraph.  Two functions that agree are
integrable together, so $F$ is Lebesgue integrable on $\mathbb{R}^n$.  Finally, a function
vanishing (almost) everywhere off a set has the same integral over that set as over the whole
space (shared context), so
\begin{equation*}
  \int_{\Omega_\varepsilon}F\,d\lambda = \int_{\mathbb{R}^n}F\,d\lambda
  \qquad (0 < \varepsilon \leq \sigma) . \tag{1}
\end{equation*}
Taking $\varepsilon = \sigma$, which is admissible since $\sigma > 0$, gives the first
conclusion.

\emph{Step 2: the truncated integral is eventually constant, so the limit is that constant.}
Consider the function
$G(\varepsilon) = \int_{\Omega_\varepsilon}F\,d\lambda$ on $(0,\infty)$, and the filter of
neighbourhoods of $0$ within $(0,\infty)$.  Every $\varepsilon$ in the intersection of
$(0,\infty)$ with the neighbourhood $(-\infty,\sigma)$ of $0$ satisfies $0 < \varepsilon < \sigma$,
and that intersection is a member of the filter; so by $(1)$,
\[
  G(\varepsilon) = \int_{\mathbb{R}^n}F\,d\lambda
  \qquad\text{for all } \varepsilon \text{ in a member of the filter,}
\]
i.e.\ $G$ agrees with the constant function $\int_{\mathbb{R}^n}F\,d\lambda$ along that filter.

Clause (2) of \noderef{PVCoreSmooth}, applied to the same data, states that $G(\varepsilon)$
tends to $\Phi_j^\varphi(z)$ as $\varepsilon \downarrow 0$, i.e.\ along that same filter.  Two
functions agreeing along a filter have the same limits, so the constant function
$\int_{\mathbb{R}^n}F\,d\lambda$ also tends to $\Phi_j^\varphi(z)$ along it.  But a constant
function tends to its own value along every filter, and the filter of neighbourhoods of $0$
within $(0,\infty)$ is not the trivial filter (every such neighbourhood contains points of
$(0,\infty)$, since $0$ is a limit point of $(0,\infty)$), so limits along it are unique.
Therefore
\[
  \Phi_j^\varphi(z) = \int_{\mathbb{R}^n}F\,d\lambda ,
\]
which is the second conclusion.
\end{proof}
